\section{Two raw controlled realizations}\label{sec:realizations}
The following proofs verify the mother theorem independently. They do not infer the general statement from a model-specific rate. Their density and accessibility mechanisms were developed in the retained controlled and finite-acquisition articles \cite{GTFControlled,GTFFinite}; we give the verifications here because the stopped law uses every hypothesis, including the preparation domain.

\begin{theorem}[Controlled Gaussian hidden-state experiments]\label{thm:gaussian}
Let the hidden physical state have $q\ge2$ values. Let $P$ be an invertible column-stochastic matrix with $P_{ij}\ge a_0>0$. Before the transition, a chosen sensor $a$ reports a Gaussian vector with mean $m_i^a$ conditional on hidden state $i$, covariance $\Sigma_a>0$, and report dimension $q-1$. The finite sensor set has full-rank likelihood-difference matrices. The hidden state then transitions by $P$. A protected categorical audit of the current state is scored by squared probability loss. All conclusions of Theorems~\ref{thm:main} and \ref{thm:cap}, and Corollary~\ref{cor:internal}, hold with $d=q-1$. The comparison includes every full-history adaptive sensor policy, not just a fixed sensor word.
\end{theorem}
\begin{proof}
The categorical audit separates the probability vector. Every legal future instrument-word probability is linear in this vector, yielding the finite predictive quotient and its positive simplex cone directly from the raw hidden-state kernel. All simplex preparations and their mixtures are allowed. Relative to Lebesgue measure the raw instrument is
\[
 L_a(y)u=P\diag(f_1^a(y),\ldots,f_q^a(y))u.
\]
A positive normalized mixture of the finitely many Gaussian densities can instead be used as reference probability, with the corresponding likelihood ratio. This changes no acquired law. The order is report first, transition second.

Take $K=P\Delta_{q-1}$ and $D=\{u\in\Delta_{q-1}:u_i\ge a_0/2\}$. Relative to the simplex hyperplane, $K$ is full-dimensional, compact and contained in the interior of $D$. All raw outputs are in $K$, and each likelihood is positive on $D$. Linearity gives report total-variation continuity; dominated integration gives weak row continuity. The loss is $g(u)=1-\norm u^2$ on the probability hyperplane, with its actual categorical variance.

Before applying $P$, let $w$ be the normalized Bayes vector and $z_i=\log(w_i/w_q)$. Then
\[
 z_i=\log(u_i/u_q)+(m_i^a-m_q^a)^T\Sigma_a^{-1}y
 -\tfrac12\bigl[(m_i^a)^T\Sigma_a^{-1}m_i^a-(m_q^a)^T\Sigma_a^{-1}m_q^a\bigr].
\]
Under the true report law this is a finite mixture of nonsingular Gaussian vectors. Compact $D$, finitely many actions, positive covariance eigenvalue bounds and the nonzero likelihood-rank determinants give a uniform density bound $C\exp(-c\norm z^2)$. The inverse logistic Jacobian is a fixed multiple of $(\prod_iw_i)^{-1}$, bounded by $C\exp(C\norm z)$. Thus Gaussian decay dominates the Jacobian on the entire simplex, including near every boundary face. The invertibility of $P$ preserves a uniform upper density on $K$.

On a fixed smaller interior simplex $W_*$, the same Gaussian-mixture formula gives a uniform strictly positive density lower for every $u\in D$ and sensor. Its image $PW_*$ supports a common normalized-volume measure $\zeta$ with $\mu_{a,u}\ge\epsilon\zeta$. Hence \eqref{eq:overlap} holds with $b=1$ even for adaptive policies. Relative hyperplane grids have null boundaries. This proves every certificate; no tail, unsuccessful report or acquisition cost was discarded. Table means are programmed constants for the known preparation and hazard, not a free exact filter register.
\end{proof}

\begin{theorem}[Noncommuting controlled quantum experiments]\label{thm:quantum}
Let $D_q\ge2$, $0<\eta,\lambda<1$ and a finite nonempty family of unitaries $U_a$ be fixed. On the unit ball of traceless Hermitian matrices, with normalized Hilbert--Schmidt volume $\lambda_0$, put $E_H=I+\eta H$, $K_H=\sqrt{E_H}$, $\tau_*=I/D_q$, and
\begin{equation}\label{eq:quantum}
 L_a(H)\sigma=\lambda U_aK_H\sigma K_HU_a^*
           +(1-\lambda)\tr(E_H\sigma)\tau_*.
\end{equation}
The action is chosen before $H$. Use a fixed informationally complete finite positive-operator-valued measurement for the protected categorical audit. Theorems~\ref{thm:main} and \ref{thm:cap}, and Corollary~\ref{cor:internal}, hold in its induced affine probability chart with $d=D_q^2-1$. The observer has classical labels only; no quantum memory is provided. Constants may deteriorate near degenerate endpoints of $\eta,\lambda$.
\end{theorem}
\begin{proof}
The effects integrate to $I$ and satisfy $(1-\eta)I\le E_H\le(1+\eta)I$. Thus \eqref{eq:quantum} is a completely positive normalized instrument: an unsharp measurement, a controlled unitary, and depolarization. Its report density $1+\eta\tr(H\sigma)$ depends on the incoming state. This is a repeated-measurement instrument in the operational sense \cite{DL1970}; the unsharp effects generally do not commute.

Informational completeness separates density operators. Every future instrument-word expectation is affine in the density operator. The physical density base and positive cone therefore give the stipulated predictive realization and a Borel preparation section. The audit chart is injective affine; its Euclidean Brier metric is equivalent, with fixed constants, to Hilbert--Schmidt distance. The actual sampling variance is not removed.

Let $K=\{\sigma\ge(1-\lambda)\tau_*:\tr\sigma=1\}$ and $D=\{\sigma\ge(1-\lambda)\tau_*/2:\tr\sigma=1\}$. All normalized raw outputs lie in $K$, even from a rank-deficient physical input. There is a compact relative buffer. Positivity, linearity and the compact effect domain give all domination and report-continuity conditions.

For fixed interior $\rho$, the map $E\mapsto\sqrt E\rho\sqrt E/\tr(E\rho)$ on trace-$D_q$ positive effects has inverse
\[
 B_{\rho,\sigma}=\rho^{-1/2}(\rho^{1/2}\sigma\rho^{1/2})^{1/2}\rho^{-1/2},
 \qquad E=\frac{D_qB_{\rho,\sigma}^{\,2}}{\tr B_{\rho,\sigma}^{\,2}}.
\]
Indeed the positive solution of $B\rho B=\sigma$ is unique; substituting it proves both inverse identities. On $\rho\in D$ and the effect ball, eigenvalues have positive lower bounds, so the inverse derivative is uniformly bounded. The forward Jacobian is consequently bounded away from zero. The final depolarizing/unitary map is invertible on the trace-one hyperplane. The bounded actual report density then gives \eqref{eq:density} in dimension $D_q^2-1$.

At $\rho=\tau_*$ the output is $\tau_*+(\lambda\eta/D_q)U_aHU_a^*$. An interior effect ball covers a common neighborhood of $\tau_*$. Smooth inversion and finitely many actions imply a local actual overlap $\mu_{a,q}\ge\epsilon_0\zeta$ for every $q$ in a sufficiently small closed neighborhood $G$ of $\tau_*$.

This local overlap is reached under every adaptive policy. With $W(\rho)=\norm{\rho-\tau_*}_{\HS}$, the matrix square-root bound and likelihood floor give
\[
 W(F_a(H,\rho))\le\lambda W(\rho)+C\lambda\eta t
 \quad(\norm H_{\HS}\le t,\ \rho\in K).
\]
One obtains the estimate by expanding numerator and denominator, using $\norm{\sqrt{I+\eta H}-I}_{\rm op}\le C\eta t$ and $\tr(E_H\rho)\ge1-\eta t$. A quiet report ball has actual conditional probability at least $(1-\eta)t^d$, for every incoming state and action. Choose $t$ so that the drift floor is smaller than one quarter of the radius of $G$, and choose $k$ so that $\lambda^k\sup_KW$ is smaller than another quarter. After the first report has entered $K$, $k$ consecutive quiet reports bring the exact predictive state into $G$, regardless of all feedback choices. The next row contributes the local overlap. Its genuine subprobability weight is at least $\epsilon_0[(1-\eta)t^d]^k$. Thus \eqref{eq:overlap} holds with $b=k+2$. Every report in this argument is paid; restricting to a report word is only a mass witness, not a postselected new experiment. The general stopped theorem now applies. No infinite-dimensional operator closure is used.
\end{proof}

\begin{corollary}[Raw changing acquired rank]\label{cor:rank}
After the instrument in Theorem~\ref{thm:quantum}, independently select a visible mark $j$ with probability $w_j$ and apply a trace-preserving unital conditional expectation $\mathcal E_j$ onto a matrix subalgebra. Include the identity with positive weight; allow complete depolarization as an atom. Then Theorem~\ref{thm:strata} applies, with $d_j$ the affine dimension of the trace-one part of that subalgebra.
\end{corollary}
\begin{proof}
The raw marked instrument is $w_j\mathcal E_jL_a(H)$. It is positive and normalized after summing marks. Each expectation fixes $\tau_*$ and preserves the positive spectral floor. In Hilbert--Schmidt coordinates it is an orthogonal projection. Fubini's theorem projects a bounded density on compact $K$ to a bounded relative density; the bound is the original bound times a maximal compact-fiber volume. A fixed audit chart changes only constants. Complete depolarization gives the exact atom $\tau_*$.

For feedback overlap, restrict the quiet report word and its final overlapping report to the identity mark. Multiply the genuine mass witness by $w_{\rm id}^{k+1}$. All other marks continue to execute on the complementary event with their unchanged weights. Relative grids provide the required cells. This is a change of acquired predictive rank, for example between full density coordinates and diagonal coordinates. The physical matrices stay strictly positive; it is not a claim of eigenvalue-rank change or a theorem about arbitrary pure-state boundaries.
\end{proof}
